\begin{rSection}{Research Experience}

\begin{rSubsection}{Keysight Technologies}{May 2026 -- Aug 2026}{Research Intern --- Machine Learning Security (Applied ML \& Fault Injection)}{Santa Rosa, CA}
\item Researching model-guided optimization and adaptive-learning methods that automatically tune fault-injection parameters to improve vulnerability-detection coverage over baseline sweep-based flows.
\item Designing experiments and telemetry instrumentation, scalable hyperparameter search, and reliability metrics; delivering a prototype demonstration on production device-security test tooling.
\end{rSubsection}

\begin{rSubsection}{The University of Kansas}{Jan 2022 -- Present}{Graduate Research Assistant, Hardware Security (Advisor: Dr.\ Tamzidul Hoque)}{Lawrence, KS}
\item \textbf{Hardware/architecture security:} designed Trojan-assisted meta-attacks that systematically locate and neutralize design-for-security primitives, demonstrating authentication bypass of ring-oscillator PUFs and unlocking of dynamically obfuscated scan chains (ACM TODAES 2026).
\item \textbf{Applied cryptography:} led a security analysis of lightweight homomorphic obfuscation schemes that protect software secrets against hardware Trojans, characterizing their guarantees and limits (IACR TCHES 2026).
\item \textbf{Software/PL-layer protection:} built HOACS, a compiler-based (LLVM) and C-library framework that keeps secrets encoded during execution on untrusted COTS processors via residue number coding --- protecting cryptographic and arithmetic workloads without hardware changes.
\item \textbf{ML for security:} developed golden-chip-free, ML-assisted power side-channel methods for detecting and localizing hardware Trojans on FPGAs (Journal of Hardware and Systems Security 2026; IEEE NAECON 2023).
\item \textbf{Tooling \& experimental R\&D:} built SPHINX, an automated framework that extracts security primitives from gate-level netlists; proposed a gray-box, pre-deployment evaluation methodology for opaque COTS hardware using power/EM side channels; demonstrated power side-channel and fault attacks on the ASCON standard and evaluated countermeasures.
\item \textbf{Collaboration \& communication:} co-authored with teams at KU, the University of Florida, and Wright State University; mentored four NSF REU undergraduates across two institutions; co-taught a hands-on hardware course and delivered side-channel tutorials and research talks.
\end{rSubsection}

\end{rSection}
